\section{Spectral resolution: two exact tests for a proposed BGS argument}
\label{sec:spectral-resolution}
The shell theorems determine finite-rank analytic data exactly, but not with
uniform microscopic resolution. This section returns to propagator spectra.
It gives necessary safeguards for a universality proof, not a proof of BGS.

\begin{proposition}[Simple spectra invisible to a long initial trace window]
For $N=2M$, let $U_N$ have the $N$th roots of unity as eigenvalues and let $V_N$
have the $2M$ distinct eigenvalues
\[
 e^{2\pi i j/M+i\pi/(3M)},\quad e^{2\pi i j/M-i\pi/(3M)},
 \qquad 0\le j<M.
\]
Then $\operatorname{Tr}U_N^k=\operatorname{Tr}V_N^k=0$ for every
$1\le k<M$. Nevertheless, with global mean-spacing unfolding, their cyclic
nearest-neighbour spacing measures are exactly
\[
 \mu_{U_N}=\delta_1,\qquad
 \mu_{V_N}=\tfrac12\delta_{2/3}+\tfrac12\delta_{4/3}.
\]
Their Wasserstein distance of order one is $1/3$, independently of $N$.
For their spectral determinants $D_U(z)=\det(I-zU_N)$ and $D_V(z)=\det(I-zV_N)$,
\[
 D_U(z)=1-z^{2M},\qquad D_V(z)=1-z^M+z^{2M}.
\]
In particular $\sup_{|z|=r}|D_U(z)-D_V(z)|\le3r^M$ for $0<r<1$.
\end{proposition}
\begin{proof}
Summing an $M$-term geometric progression gives the trace assertion and,
more generally, $N^{-1}\operatorname{Tr}V_N^{kM}=\cos(k\pi/3)$.
The two phases in each cell have unfolded distance $2/3$, and the next gap
has length $4/3$. Their determinants are the products of
$1-e^{i\pi/3}z^M$ and $1-e^{-i\pi/3}z^M$. The claimed difference bound
follows from $D_V-D_U=-z^M+2z^{2M}$. The distance to the point mass at one
is the mean absolute deviation, namely $1/3$.
\end{proof}

For the normalized Carath\'eodory transform
\[
 G_U(z)=N^{-1}\operatorname{Tr}\{(I+zU)(I-zU)^{-1}\}
       =1+2\sum_{k\ge1}z^kN^{-1}\operatorname{Tr}U^k,
\]
the same trace formula gives
\[
 \sup_{|z|=r}|G_{U_N}(z)-G_{V_N}(z)|
       \le\frac{4r^M}{1-r^M}.
\]
Consequently the determinants approach each other and these transforms both
approach one on the circles $r_N=\exp(-s_N/N)$ whenever $s_N\to\infty$.
Their microscopic spacing measures remain separated. This is an explicit
resolution obstruction even with simple spectra and bounded unfolded gaps.
It does not say that either clock construction is a quantization of a
chaotic map. It says that short traces and unquantified mesoscopic resolvent
convergence cannot, as abstract spectral hypotheses, identify the local law.
At $r_N=\exp(-s/N)$ with fixed $s$, the quantity $r_N^M=e^{-s/2}$ no longer
vanishes; the missing boundary-layer scale is precisely of order $1/N$.

\begin{proposition}[First-order Egorov alone cannot force local repulsion]
Suppose $U_N$ is any unitary quantization satisfying, for fixed observables $a$,
\[
 \|U_N^*\operatorname{Op}_N(a)U_N-\operatorname{Op}_N(a\circ T)\|
       =O_a(N^{-1}),\qquad \sup_N\|\operatorname{Op}_N(a)\|<\infty.
\]
Fix an integer $L\ge2$ and take $M_N=\lfloor N/L\rfloor$, for $N\ge L$.
There are unitaries $\widetilde U_N$ with
$\|\widetilde U_N-U_N\|=O_L(N^{-1})$, the same displayed Egorov estimate,
and at least $N-M_N$ zero cyclic spacings. There is also a simple-spectrum
version whose limiting spacing measures still have an atom of mass at least
$1-1/L$ at zero, whenever a weak limit is taken.
\end{proposition}
\begin{proof}
Use the spectral theorem to round every eigenangle to a nearest $M_N$th
root of unity. The change in operator norm is at most
$2\sin(\pi/(2M_N))\le\pi/M_N=O_L(N^{-1})$.
There are at most $M_N$ occupied roots, hence at least $N-M_N$ zero spacings.
For any bounded operator $B$ the conjugation difference is at most
$2\|B\|\|\widetilde U_N-U_N\|$, which preserves the assumed Egorov bound.

For the simple-spectrum version, split each occupied root into distinct
angles in an arc of half-width $\varepsilon_N/N$, where
$\varepsilon_N>0$ tends to zero. These arcs can be chosen disjoint.
At least $N-M_N$ spacings are then at most
$2\varepsilon_N/N$, or $\varepsilon_N/\pi$ after unfolding by $N/(2\pi)$.
The extra norm change is at most $\varepsilon_N/N$ and preserves Egorov.
For any fixed $\eta>0$ the limiting mass in the closed interval $[0,\eta]$
is at least $1-1/L$ by the portmanteau inequality; let $\eta\downarrow0$.
\end{proof}

This conclusion applies to the stated weak quantization axioms regardless of
the classical map. It is not a counterexample to a formulation imposing a
fixed smooth microlocal symbol class or a specified canonical quantization:
the spectral rounding is not shown to belong to such a class. In particular,
neither this construction nor the clock pair can replace the missing
Heisenberg-scale joint-trace estimate for the maintained kicked propagator.
Their purpose is to reject proof attempts that use only first-order Egorov,
classical covariance positivity, control Lie rank, or coarse determinant data.
